% Chapter 15: Analogies and Their Obstructions
% Organized by OBSTRUCTION TYPE, not by analogy source.
% Content: ALL 32 analogy edges from the graph.

\chapter{Analogies and Their Obstructions}
\label{ch:analogies}

When you read that ``compression is like VC dimension'' or that ``stability
is related to PAC-Bayes,'' you are being handed an analogy without its
expiry date. The remark could mean that both quantities measure the same
underlying thing; it could mean that one bounds the other in one direction
only; or it could mean that the objects involved belong to different parts
of the type system and no formal connection will ever hold. Without knowing
which, you cannot tell whether to try proving the equivalence, whether to
expect a counterexample, or whether to search for a new bridging concept
entirely. The analogy, stated informally, is a map with the legend missing.

This chapter supplies the legend. The concept graph contains 32 edges of
type \rel{analogy}, and each is treated here as a formal object with three
components:

\begin{enumerate}[leftmargin=*, nosep]
\item A \emph{source} and \emph{target}: two concepts that share a
structural parallel.
\item A \emph{shared structure}: the precise mathematical property that
makes the analogy plausible.
\item An \emph{obstruction}: the precise reason the analogy fails to be
a theorem, classified by \emph{type}.
\end{enumerate}

The obstruction type carries more information than the analogy itself. Two
analogies can look identical at the surface level (``X is like Y because
both measure complexity'') while failing for entirely different reasons:
one because the objects live in different type systems, another because the
conjectured equivalence is plausible but still unproved. The type identifies the reason the analogy fails, and that
determines whether the gap is a research target, a structural impossibility,
or something in between.

The 32 analogy edges distribute across six obstruction types:

\begin{center}\small
\begin{tabular}{llr}
\toprule
\textbf{Obstruction type} & \textbf{Meaning} & \textbf{Count} \\
\midrule
Type mismatch & Objects live in different formal categories & 12 \\
Missing equivalence witness & Plausible but unproved & 9 \\
One-way theorem only & One direction proved, reverse open/false & 4 \\
Data model mismatch & Incompatible data-generating assumptions & 3 \\
Proof method mismatch & Same structure, non-transferable proofs & 3 \\
Success criterion mismatch & Different optimization objectives & 1 \\
\bottomrule
\end{tabular}
\end{center}

The plurality are type mismatches: the objects being compared belong to
different formal categories, and making the analogy a theorem would require
new bridging concepts that do not yet exist. The second largest group, 9
missing equivalence witnesses, is the frontier: analogies where someone
proving the right conjecture would convert an informal remark into a
theorem. The sample compression conjecture sits here. So does the
MML--MDL relationship. These are not structurally blocked; they are open.


% ================================================================
% SECTION 1: TYPE MISMATCH (12 edges)
% ================================================================

\section{Type Mismatch}
\label{sec:type-mismatch}

A type mismatch arises when the source and target of an analogy belong
to different formal categories in the concept graph. The structural
parallel is real: both concepts do something similar. But they operate on
different mathematical objects, and no bridging theorem connects the two
type systems.

These are the most common obstruction, accounting for 12 of the 32 edges,
and the least likely to be resolved. The gap is not waiting for a clever
proof; it is a mismatch in the formalism itself. The analogy is not a
conjecture. It is a description of two things that happen to resemble each
other across a category boundary.

\begin{table}[ht]
\centering\small
\caption{Twelve analogies that fail because the source and target live in
different formal categories. The ``type gap'' column names the specific
category boundary that blocks a formal theorem. For each pair, the shared
structure is real; the gap is why that shared structure cannot be stated
as an equivalence without a bridging concept that does not yet exist.}
\label{tab:type-mismatch}
\begin{tabularx}{\textwidth}{@{}p{4.2cm}p{3.2cm}X@{}}
\toprule
\textbf{Analogy} & \textbf{Shared structure} & \textbf{Type gap} \\
\midrule

\nde{algorithmic\_stability} $\leftrightarrow$ \nde{pac\_bayes\_bound}
& Both bound generalization gap
& Stability is algorithm-dependent (perturbation of $A$). PAC-Bayes is
posterior-dependent (divergence $\KL(Q \| P)$). An algorithm defines a
posterior, but the primary objects differ. \\[6pt]

\nde{information\_theoretic\_bound} $\leftrightarrow$
\nde{pac\_bayes\_bound}
& Both use information-divergence measures
& Info-theoretic bounds use $I(W; S)$ (mutual information between
algorithm output and training set). PAC-Bayes uses $\KL(Q \| P)$
(posterior vs.\ prior). Different information measures on different
objects. Hellstr\"om et al.\ (2023) partially unify via change-of-measure,
but the objects remain distinct. \\[6pt]

\nde{fundamental\_theorem} $\leftrightarrow$
\nde{algorithmic\_stability}
& Both characterize learnability
& The fundamental theorem is a characterization (equivalence result).
Stability is a complexity measure. The structural parallel exists but the
formal types are incompatible without a bridging theorem. \\[6pt]

\nde{sq\_dimension} $\leftrightarrow$ \nde{vc\_dimension}
& Both measure concept class complexity
& SQ dimension measures pairwise correlation structure under a
distribution. VC dimension measures distribution-free shattering.
Exponential gap possible (parities: $\VC = n$, $\SQdim = 2^n$). \\[6pt]

\nde{grammar\_induction} $\leftrightarrow$ \nde{gold\_theorem}
& Both concern language learning
& Grammar induction is a process; Gold's theorem is an impossibility
result. Different concept categories. \\[6pt]

\nde{concept\_drift} $\leftrightarrow$ \nde{ex\_under\_drift}
& Both concern non-stationary learning
& Concept drift is a process; Ex-under-drift is a success criterion.
Different concept categories. \\[6pt]

\nde{lifelong\_learning} $\leftrightarrow$ \nde{meta\_learner}
& Both concern learning across tasks
& Lifelong learning is a process; meta-learner is a learner type.
Different concept categories. \\[6pt]

\nde{background\_knowledge} $\leftrightarrow$ \nde{advice}
& Both provide side information to a learner
& Background knowledge is a process; advice is a data presentation.
Different concept categories. \\[6pt]

\nde{mdl} $\leftrightarrow$ \nde{kolmogorov\_complexity}
& Both measure description length
& MDL uses finite computable codes. Kolmogorov complexity uses the
shortest program on a universal Turing machine (generally
incomputable). The parallel is ``short description = good model,''
but MDL's codes are constructive while Kolmogorov's are not. \\[6pt]

\nde{srm} $\leftrightarrow$ \nde{inductive\_bias}
& Both constrain hypothesis selection
& SRM is a model selection principle; inductive bias is a base type.
The structural parallel exists but the categories are incompatible. \\[6pt]

\nde{bayesian\_inference} $\leftrightarrow$ \nde{inductive\_bias}
& Both encode prior assumptions
& Bayesian inference is a model selection method (probability over
hypotheses). Inductive bias is a base type (constraint on learner
behavior). The prior \emph{is} a specific form of bias, but the type
systems differ. \\[6pt]

\nde{occam\_algorithm} $\leftrightarrow$ \nde{compression\_scheme}
& Both embody ``short representation $\Rightarrow$ generalization''
& Occam uses description length of hypotheses (syntactic: bits to
encode $h$). Compression uses a subset of training data (extensional:
examples to reconstruct $h$). The notion of shortness differs. \\

\bottomrule
\end{tabularx}
\end{table}

In each of these 12 cases, the analogy identifies something real: a
structural parallel that is not an accident. But the source and target live
in different parts of the type system, and formalizing the parallel requires
a bridging theorem relating the two types. No impossibility result rules out
such a bridge for any individual pair. One pair, information-theoretic
bounds versus PAC-Bayes, already has a partial bridge via change-of-measure
\cite{HellstromDurisi2020}. The table records where bridges are currently
absent, not where they are provably out of reach.


% ================================================================
% SECTION 2: MISSING EQUIVALENCE WITNESS (9 edges)
% ================================================================

\section{Missing Equivalence Witness}
\label{sec:missing-equiv}

A missing equivalence witness marks an analogy where the formal equivalence
looks plausible but has no proof. These are the active frontier of the
field: the places where a theorem is conceivably available, and where
finding one would matter.

Not all nine edges here are genuinely open. Three are settled and appear in
this section only because the analogy is commonly stated as if it were an
equivalence. The mistake-bound/Littlestone pair \emph{is} an equivalence
(Littlestone's theorem, proved below). The Rademacher/VC and
label-complexity/sample-complexity pairs are established
\emph{non}-equivalences: the separations are known, not conjectural.
The genuinely open conjectures are compression/VC and the two MML edges.
The distinction between ``open conjecture'' and ``known separation carrying
the same color'' is load-bearing; the exercises at the end of the chapter
press directly on it.

\begin{obstbox}
\textbf{Compression $\leftrightarrow$ VC dimension.}

\medskip\noindent
The sample compression conjecture asserts that every concept class with
VC dimension $d$ admits a sample compression scheme of size $O(d)$.
Moran and Yehudayoff~\cite{MoranYehudayoff2016} proved that finite VC
dimension implies \emph{some} finite compression scheme (of size
exponential in $d$), establishing one
direction.\formal{FLT_Proofs.Theorem.PAC.fundamental_vc_compression} The conjecture that
compression size is $O(d)$ remains open. If proved, this would upgrade
the analogy to an equivalence; if disproved, it would become a
separation. It is the most important open problem connected to the
analogy edges in the graph.
\end{obstbox}

\begin{obstbox}
\textbf{Rademacher complexity $\leftrightarrow$ VC dimension.}

\medskip\noindent
The two are asymptotically related: $\Rad_n(\mathcal{H}) \leq
O\!\left(\sqrt{\VC(\mathcal{H})/n}\right)$. Both measure the complexity
of a hypothesis class, and for binary classification both control
generalization. But Rademacher complexity is data-dependent
(computed from a specific sample and distribution), while VC dimension is
distribution-free. No equivalence holds: Rademacher complexity can be
much smaller than the VC bound for benign distributions, and the two
quantities can disagree on the ranking of hypothesis classes.
\end{obstbox}

\begin{obstbox}
\textbf{VC characterization $\leftrightarrow$ Littlestone
characterization.}

\medskip\noindent
Both theorems have the same logical form: ``a class is learnable in
paradigm $P$ if and only if combinatorial dimension $d$ is finite.''
The parallel is exact at the level of logical structure. But the
dimensions measure different adversarial strengths (batch (VC) vs.\
adaptive (Littlestone)), and $\VC \leq \Ldim$ with no reverse
bound.\formal{FLT_Proofs.Theorem.Separation.vcdim_le_of_mistake_bounded}
The gap can be infinite (thresholds on $\R$). Alon et
al.~\cite{AlonLivniMalliarisMoran2019} connected the two via
differential privacy, but no equivalence between the dimensions exists.
\end{obstbox}

\begin{obstbox}
\textbf{MML $\leftrightarrow$ MDL.}

\medskip\noindent
Minimum Message Length (Wallace--Freeman) and Minimum Description Length
(Rissanen) both select hypotheses by minimizing a two-part code: the
model plus the data given the model. They are often treated as variants
of the same idea. However, MML uses a Bayesian prior explicitly and
integrates over the parameter space, while MDL uses a minimax-optimal
universal code. No formal theorem establishes their equivalence or
separation; the two traditions developed independently and use different
mathematical frameworks.
\end{obstbox}

\begin{obstbox}
\textbf{MML $\leftrightarrow$ Bayesian inference.}

\medskip\noindent
MML selects the hypothesis that minimizes total message length; Bayesian
inference selects (or averages over) hypotheses according to posterior
probability. For certain prior--likelihood pairs, MML and MAP
(maximum a posteriori) estimates coincide asymptotically. But MML is a
coding-theoretic criterion and Bayesian inference is a probabilistic
criterion; the equivalence is approximate and model-dependent, holding for
particular prior--likelihood pairs rather than as a general theorem.
\end{obstbox}

\begin{obstbox}
\textbf{PAC learning $\leftrightarrow$ posterior consistency.}

\medskip\noindent
Both are success criteria requiring convergence to the truth. PAC
demands uniform convergence in probability over all distributions;
Bayesian consistency demands that the posterior concentrates on the true
parameter. The structural parallel is ``learning = convergence,'' but
PAC convergence is frequentist and distribution-free, while posterior
consistency is Bayesian and prior-dependent. No general theorem
connects the two: a class can be PAC learnable but Bayesian-inconsistent
under a misspecified prior, and vice versa.
\end{obstbox}

\begin{obstbox}
\textbf{Mistake bound $\leftrightarrow$ Littlestone dimension.}

\medskip\noindent
The optimal mistake bound for a concept class equals its Littlestone
dimension (this is Littlestone's
theorem).\formal{FLT_Proofs.Theorem.Online.optimal_mistake_bound_eq_ldim}
However, the \emph{analogy} edge in the graph refers to the structural
parallel between the two as complexity measures in their own right. The mistake bound is defined
operationally (worst-case mistakes of the best algorithm); the
Littlestone dimension is defined combinatorially (depth of the largest
shattered tree). Their equivalence is a theorem, but as stand-alone
objects they measure different things, and extending either to new
settings (e.g., partial monitoring, bandit feedback) produces divergent
generalizations.
\end{obstbox}

\begin{obstbox}
\textbf{Label complexity $\leftrightarrow$ sample complexity.}

\medskip\noindent
Both count the amount of data needed for learning, but label complexity
(from active learning) counts the number of \emph{labels queried},
while sample complexity counts the number of \emph{labeled examples
drawn}. In active learning, the learner sees many unlabeled points but
requests labels selectively, so label complexity can be exponentially
smaller than sample complexity. No equivalence exists; the gap is the
content of active learning theory.
\end{obstbox}

\begin{obstbox}
\textbf{Lifelong learning $\leftrightarrow$ concept drift.}

\medskip\noindent
Both concern learning in non-stationary environments. Lifelong learning
emphasizes knowledge transfer across a sequence of tasks; concept drift
emphasizes adaptation to a changing target within a single task. The
structural parallel is ``the learning problem changes over time,'' but
the formalization differs: lifelong learning uses a task distribution,
concept drift uses a time-varying concept. No formal theorem connects
the two frameworks.
\end{obstbox}


% ================================================================
% SECTION 3: ONE-WAY THEOREM ONLY (4 edges)
% ================================================================

\section{One-Way Theorem Only}
\label{sec:one-way}

In these four analogies, one direction has been proved and the other has
not. The structural parallel is real, but it runs one way. What is at stake
here is whether the learner conflates a proved implication with a
biconditional: assuming symmetry that the mathematics does not support, and
then treating the converse as obvious when it is either open or false.

\begin{obstbox}
\textbf{Sample complexity $\leftrightarrow$ VC dimension}
(a theorem that an identity-reading would erase).

\medskip\noindent
The fundamental theorem establishes the biconditional
$\VC(\mathcal{H}) < \infty \Leftrightarrow$ finite sample
complexity.\formal{FLT_Proofs.Theorem.PAC.fundamental_theorem} Unlike the
other entries in this section, this relationship runs \emph{both} ways; it
appears here as a cautionary case.  The analogy is usually stated as an
identity (``sample complexity just \emph{is} VC dimension''), which
collapses a two-directional theorem (an explicit upper bound one way, a
matching lower bound the other, proved by different arguments) into a
definition, erasing exactly the content the fundamental theorem supplies.
\end{obstbox}

\begin{obstbox}
\textbf{Compression scheme $\rightarrow$ sample complexity} (but not
$\leftarrow$).

\medskip\noindent
A compression scheme of size $k$ implies PAC sample complexity
$O(k \cdot \log m / \varepsilon)$ (Littlestone--Warmuth, Moran--Yehudayoff).
But no theorem says that optimal sample complexity determines optimal
compression size. A class might be learnable with few samples yet require
a large compression scheme. The compression direction is proved;
the reverse is open.
\end{obstbox}

\begin{obstbox}
\textbf{Covering number $\rightarrow$ VC dimension} (but not
$\leftarrow$).

\medskip\noindent
The Sauer--Shelah lemma gives: $\log N(\varepsilon, \mathcal{H},
L_1(P)) \leq d \cdot \log(2e/\varepsilon)$, bounding covering numbers
from VC dimension.\formal{FLT_Proofs.Complexity.Rademacher.sauer_shelah_exp_bound} But covering numbers are defined for arbitrary
metric spaces, while VC dimension applies only to binary hypothesis
classes. Finite VC dimension suffices for covering number
bounds; it is not required for them, since the covering number framework is
strictly more general and bounds spaces that have no VC dimension at all.
\end{obstbox}

\begin{obstbox}
\textbf{Covering number $\rightarrow$ Rademacher complexity} (but not
$\leftarrow$).

\medskip\noindent
Dudley's entropy integral bounds Rademacher complexity from above using
covering numbers: $\Rad_n(\mathcal{H}) \leq \inf_{\alpha > 0}
\left[ 4\alpha + \frac{12}{\sqrt{n}} \int_\alpha^1
\sqrt{\log N(\varepsilon, \mathcal{H}, L_2)} \, d\varepsilon \right]$.
This bound can be loose; tighter chaining methods (Talagrand's generic
chaining) and direct Rademacher estimates sometimes give better results.
The Dudley bound is one-directional: covering numbers control
Rademacher, but Rademacher complexity does not determine covering
numbers.
\end{obstbox}


% ================================================================
% SECTION 4: DATA MODEL MISMATCH (3 edges)
% ================================================================

\section{Data Model Mismatch}
\label{sec:data-model}

These three analogies fail at the level of assumptions about how data
arrives. The structural parallel is real: the results look similar and the
intuitions transfer. But the data-generating processes are incompatible.
A result proved under i.i.d.\ sampling does not migrate to adversarial
sequences, and a result that assumes stationarity cannot be directly applied
where non-stationarity is the explicit object of study. The analogy works
as intuition; it breaks as a formal argument the moment the data model is
stated precisely.

\begin{obstbox}
\textbf{Concept drift $\leftrightarrow$ online learning.}

\medskip\noindent
Both are sequential learning settings, but concept drift uses i.i.d.\
draws per time step from a slowly changing distribution, while online
learning uses adversarially chosen instances from a fixed concept class.
The shared structure is ``sequential prediction under non-stationarity,''
but the non-stationarity is stochastic in one case and adversarial in
the other. Results from one setting do not transfer: online regret
bounds assume a fixed comparator class, while drift bounds assume a
rate of distributional change.
\end{obstbox}

\begin{obstbox}
\textbf{Advice reduction $\leftrightarrow$ sample complexity.}

\medskip\noindent
Both quantify ``how much side information shrinks the learner's search
space.'' But advice operates in the Gold model (infinite stream,
computable learner), while sample complexity operates in the PAC model
(finite i.i.d.\ sample, distribution-free guarantees). No formal theorem
bridges the two frameworks. The structural parallel (``side information
helps'') is too vague to formalize without choosing a specific data
model, and the two models are incompatible.

\hfill\cite{KaufmannStephan2001}
\end{obstbox}

\begin{obstbox}
\textbf{Granger causality $\leftrightarrow$ concept drift.}

\medskip\noindent
Both operate on time-series data, but Granger causality assumes
stationarity within estimation windows (testing whether past values
of $X$ improve prediction of $Y$), while concept drift explicitly
models non-stationarity. The structural parallel is ``temporal
dependence in sequential data,'' but the assumptions about the data
process are contradictory: Granger requires stationarity; drift requires
its absence.
\end{obstbox}


% ================================================================
% SECTION 5: PROOF METHOD MISMATCH (3 edges)
% ================================================================

\section{Proof Method Mismatch}
\label{sec:proof-method}

These three analogies have identical theorem structure. The result shapes
are parallel and the logical form is the same. But the proofs inhabit
different mathematical worlds: concentration inequalities in one case,
change-of-measure arguments in another, Ramsey theory in a third. The
analogy holds at the level of what is proved; the obstruction is that
proving one tells you nothing about how to prove the other. Knowing that
stability gives a generalization bound does not, by itself, give you a
proof technique for information-theoretic bounds, even though the conclusion
has the same form.

\begin{obstbox}
\textbf{Information-theoretic bound $\leftrightarrow$ algorithmic
stability.}

\medskip\noindent
Both bound the generalization gap via properties of the learning
algorithm. Stability uses sup-norm perturbation sensitivity (how much
does the output change when one training point is replaced?).
Information-theoretic bounds use mutual information $I(W; S)$ (how much
does the algorithm's output depend on the training set?). These are
connected via the chain: low stability $\Rightarrow$ low $I(W; S)$
$\Rightarrow$ low generalization gap. But the proof methods are
fundamentally different: stability arguments use coupling and
McDiarmid-type concentration; information-theoretic arguments use
change-of-measure and data processing inequalities. Neither proof
technique subsumes the other.
\end{obstbox}

\begin{obstbox}
\textbf{Ordinal VC dimension $\leftrightarrow$ mind-change ordinal.}

\medskip\noindent
Both use ordinal-valued measures to characterize learnability in
transfinite settings. Ordinal VC dimension (Martin--Sharma--Stephan)
characterizes predictive complexity; mind-change ordinals characterize
identification complexity in the Gold model. Both use ordinals, but they
measure different things. The connection is indirect, routed through predictive
complexity: ordinal VC dimension governs prediction error
bounds, while mind-change ordinals govern convergence to a correct
hypothesis. The proof techniques (Ramsey-theoretic for ordinal VC,
topological for mind-change) do not transfer.
\end{obstbox}

\begin{obstbox}
\textbf{Occam's razor $\leftrightarrow$ MDL.}

\medskip\noindent
Both assert that simpler models generalize better. Occam's razor (in the
PAC sense of Blumer et al.) is a worst-case generalization theorem:
if the hypothesis is short, its error is low with high probability.
MDL is an information-theoretic model selection principle: minimize the
total description length of model plus data-given-model. The theorem
\emph{structure} is parallel (``short $\Rightarrow$ good''), but Occam
gives frequentist, distribution-free guarantees while MDL gives
Bayesian-flavored, coding-theoretic guarantees. The proof methods
inhabit different mathematical worlds.

\hfill\cite{BlumerEhrenfeuchtHausslerWarmuth1987}
\end{obstbox}


% ================================================================
% SECTION 6: SUCCESS CRITERION MISMATCH (1 edge)
% ================================================================

\section{Success Criterion Mismatch}
\label{sec:criterion-mismatch}

A single analogy edge in the graph carries this obstruction. The type does not
reduce to the others: the two concepts share a setting and a procedure, but
they answer different questions about that setting.

\begin{obstbox}
\textbf{Granger causality $\leftrightarrow$ online learning.}

\medskip\noindent
Granger causality tests causal structure: does the past of time series
$X$ improve prediction of time series $Y$ beyond what $Y$'s own past
provides? Online learning minimizes cumulative prediction error against
an adversary. Both involve sequential prediction, but they optimize
different objectives. Granger seeks a \emph{structural} conclusion
(causal influence exists or does not); online learning seeks a
\emph{performance} guarantee (low regret). The success criteria are
fundamentally different: statistical significance of a causal test vs.\
sub-linear regret growth. No formal theorem connects them because they
answer different questions about the same sequential data.
\end{obstbox}


% ================================================================
% SECTION 7: THE OBSTRUCTION MAP
% ================================================================

\section{The Obstruction Map}
\label{sec:obstruction-map}

\Cref{fig:obstruction-map} puts all 32 analogy edges in one place, colored
by obstruction type, with a filled dot marking the five edges where
formalization has already produced a verified bridge. Reading the color
distribution is reading the map of what formal learning theory does not yet
know how to unify, and the five dots mark exactly where that effort has
succeeded so far.

\begin{figure}[ht]
\centering
\begin{tikzpicture}[
  x=1.5cm, y=1.2cm,
  concept/.style={
    rectangle, rounded corners=2pt, draw, thick,
    fill=layergray, font=\scriptsize\sffamily,
    minimum width=1.3cm, minimum height=0.5cm,
    align=center, inner sep=2pt
  },
  tm/.style={-, edgeorange, thick},        % type mismatch
  mew/.style={-, defblue, thick, dashed},   % missing equiv witness
  owt/.style={-, sepgreen, thick, dotted},  % one-way theorem
  dm/.style={-, thmred, thick, dashdotted}, % data model mismatch
  pm/.style={-, analogypurple, thick, densely dashed}, % proof method
  sc/.style={-, boundbrown, thick, densely dotted},    % success criterion
  ev/.style={circle, fill=black, draw=white, line width=0.25pt,
             inner sep=0pt, minimum size=3.6pt}, % verified-handle dot
]

% Layout by connected component: the analogy graph is disconnected, so each
% component sits in its own region (zero inter-region crossings) and is drawn
% planar. Color/style = obstruction type; a filled dot = a verified sorry-free
% kernel handle backs that edge.

% --- Complexity core (VC-hub component) ---
\node[concept] (vc)      at (1.5,3.5)  {VC dim};
\node[concept] (sq)      at (0.4,3.5)  {SQ dim};
\node[concept] (rad)     at (1.5,4.5)  {Rademacher};
\node[concept] (cov)     at (2.7,4.5)  {Covering\\number};
\node[concept] (comp)    at (1.5,2.5)  {Compression};
\node[concept] (sc_node) at (2.85,3.0) {Sample\\complexity};
\node[concept] (lc)      at (4.1,3.4)  {Label\\complexity};
\node[concept] (ar)      at (4.1,2.4)  {Advice\\reduction};

% --- Model-selection tail (hangs off compression via Occam) ---
\node[concept] (occam)   at (1.5,1.5)  {Occam};
\node[concept] (mdl)     at (2.7,1.5)  {MDL};
\node[concept] (kolm)    at (2.7,0.6)  {Kolmogorov};
\node[concept] (mml)     at (3.9,1.5)  {MML};
\node[concept] (bayes)   at (3.9,0.6)  {Bayesian};
\node[concept] (ib)      at (5.1,1.05) {Inductive\\bias};
\node[concept] (srm)     at (5.1,2.0)  {SRM};

% --- Verified dimension-equivalence pairs (isolated edges, both dotted) ---
\node[concept] (ldim)    at (3.9,4.6)  {Littlestone\\dim};
\node[concept] (mb)      at (5.1,4.6)  {Mistake\\bound};
\node[concept] (vcc)     at (3.5,5.5)  {VC charac.};
\node[concept] (lc_char) at (4.9,5.5)  {Ldim charac.};

% --- Sequential component (drift-hub) ---
\node[concept] (drift)   at (7.0,3.4)  {Concept\\drift};
\node[concept] (ol)      at (6.3,2.6)  {Online};
\node[concept] (gc)      at (7.7,2.6)  {Granger};
\node[concept] (exd)     at (7.7,4.2)  {Ex under\\drift};
\node[concept] (ll)      at (6.3,4.2)  {Lifelong};
\node[concept] (ml)      at (6.3,5.1)  {Meta-\\learner};

% --- Paradigm component (stability-hub) ---
\node[concept] (stab)    at (3.1,-0.4) {Stability};
\node[concept] (pb)      at (4.1,0.1)  {PAC-Bayes};
\node[concept] (itb)     at (4.1,-0.9) {Info-theoretic};
\node[concept] (ft)      at (2.0,-0.4) {Fund.\\Theorem};

% --- Remaining isolated edges ---
\node[concept] (pac)     at (0.4,1.5)  {PAC};
\node[concept] (post)    at (0.4,0.6)  {Posterior\\consistency};
\node[concept] (gi)      at (1.1,-1.7) {Grammar\\induction};
\node[concept] (gt)      at (2.3,-1.7) {Gold's\\theorem};
\node[concept] (bk)      at (3.7,-1.7) {Background\\knowledge};
\node[concept] (adv)     at (4.9,-1.7) {Advice};

% --- Type mismatch edges (orange, 12) ---
\draw[tm] (stab) -- (pb);
\draw[tm] (itb) -- (pb);
\draw[tm] (ft) -- (stab);
\draw[tm] (sq) -- (vc);
\draw[tm] (gi) -- (gt);
\draw[tm] (drift) -- (exd);
\draw[tm] (ll) -- (ml);
\draw[tm] (bk) -- (adv);
\draw[tm] (mdl) -- (kolm);
\draw[tm] (srm) -- (ib);
\draw[tm] (bayes) -- (ib);
\draw[tm] (occam) -- (comp);

% --- Missing equivalence witness edges (blue dashed, 9); 3 verified-backed ---
\draw[mew] (comp) -- (vc)      node[ev,pos=0.5]{}; % fundamental_vc_compression
\draw[mew] (rad) -- (vc);
\draw[mew] (vcc) -- (lc_char)  node[ev,pos=0.5]{}; % vcdim_le_of_mistake_bounded
\draw[mew] (mml) -- (mdl);
\draw[mew] (mml) -- (bayes);
\draw[mew] (pac) -- (post);
\draw[mew] (mb) -- (ldim)      node[ev,pos=0.5]{}; % optimal_mistake_bound_eq_ldim
\draw[mew] (lc) -- (sc_node);
\draw[mew] (ll) -- (drift);

% --- One-way theorem edges (green dotted, 4); 2 verified-backed ---
\draw[owt] (sc_node) -- (vc)  node[ev,pos=0.5]{}; % fundamental_theorem
\draw[owt] (comp) -- (sc_node);
\draw[owt] (cov) -- (vc)      node[ev,pos=0.5]{}; % sauer_shelah_exp_bound
\draw[owt] (cov) -- (rad);                        % Dudley: no kernel handle, undotted

% --- Data model mismatch edges (red dashdotted, 3) ---
\draw[dm] (drift) -- (ol);
\draw[dm] (ar) -- (sc_node);
\draw[dm] (gc) -- (drift);

% --- Proof method mismatch edges (purple, 3) ---
\draw[pm] (itb) -- (stab);
\draw[pm] (occam) -- (mdl);
% ordinal_vc_dim -> mind_change_ordinal not shown (nodes would clutter)

% --- Success criterion mismatch (brown, 1) ---
\draw[sc] (gc) -- (ol);

\end{tikzpicture}

\caption[Obstruction map: thirty-two analogies, six failure types, five
verified bridges]{Thirty-two cross-paradigm analogies, none of them a
theorem, fail for six classifiable reasons. Color and line style name the
obstruction type; a filled dot marks the five edges where formalization has
already produced a verified one-way bridge or equivalence
(\leanname{fundamental_theorem}). The argument the figure makes is that
analogy-failure is not case-by-case: it is taxonomic, and the six types
exhaust what the field currently has names for. The five dots show the
frontier where the taxonomy has been broken by a proof. Key:
\textcolor{edgeorange}{\rule[0.4ex]{1.1em}{0.9pt}} type mismatch (12);
\textcolor{defblue}{\textbf{-\,-\,-}} missing equivalence witness (9);
\textcolor{sepgreen}{\textbf{$\cdots$}} one-way theorem (4);
\textcolor{thmred}{\textbf{-$\cdot$-}} data-model mismatch (3);
\textcolor{analogypurple}{\textbf{-\,-}} proof-method mismatch (3);
\textcolor{boundbrown}{\textbf{$\cdots$}} success-criterion mismatch (1);
$\bullet$~verified bridge handle (5).
One proof-method edge (ordinal VC dimension $\leftrightarrow$ mind-change
ordinal) is omitted for clarity.}
\label{fig:obstruction-map}
\end{figure}

\subsection*{What the distribution tells us}

The obstruction type distribution reveals three structural facts about
formal learning theory.

\emph{First}, 12 of the 32 failures are type mismatches: the objects being
compared are not the same kind of thing. This is not a failure of proof
technique. The proofs have not simply not been found; the objects live in
different categories. Future unification will require new bridging concepts
that mediate between the existing type systems.

\emph{Second}, 9 missing equivalence witnesses mark the open frontier.
The sample compression conjecture (compression $\leftrightarrow$ VC
dimension) is the most prominent: Moran and Yehudayoff proved the
exponential-size direction in 2016, and the $O(d)$ conjecture remains open.
The Rademacher--VC relationship and the characterization-theorem parallel
(VC $\leftrightarrow$ Littlestone) are equally fundamental. A proof on any
of these would restructure the analogy map.

\emph{Third}, the 4 one-way theorems show where the asymmetry is built into
the objects themselves. Covering numbers bound Rademacher complexity via
Dudley's entropy integral, and covering numbers are bounded by VC dimension
via Sauer--Shelah. Both directions flow one way only. The asymmetry is not
an artifact of proof technique; it reflects a generality hierarchy: covering
numbers apply to metric spaces, Rademacher complexity to function classes,
VC dimension to binary classes. Each step downward tightens the bound and
loses applicability.

The 32 edges and their obstructions together form a map of what formal
learning theory does not yet know how to unify. Understanding why an analogy
fails is, in this field, as informative as understanding why a theorem
holds.


% ================================================================
% SECTION 8: REALIZED ANALOGIES (the verified counterpoint)
% ================================================================

\section{When the Analogy Is a Theorem}
\label{sec:realized-analogies}

Some analogies carry no obstruction at all. They began as loose parallels
and someone proved them. The cleanest is now machine-checked.

The flagship is the correspondence \emph{online learning
$\leftrightarrow$ zero-sum games}. A learner facing an adversary and a row
player in a matrix game look, at first glance, merely analogous: both react
to a worst-case opponent. The analogy is discharged by the
multiplicative-weights (Hedge) algorithm. Run on the game matrix, it
produces a distribution over rows whose guaranteed payoff is within
$\varepsilon$ of the game's value. No-regret online play is not
\emph{like} approximate game solving; it \emph{is} approximate game solving.

\begin{thm}[Multiplicative weights solve the game]
\label{thm:mwu-minimax}
Let $M \colon R \times C \to \{0,1\}$ be a finite Boolean game whose value
is at least~$v$ (every column distribution admits a row response of payoff
$\geq v$). For every $\varepsilon > 0$, the multiplicative-weights
algorithm returns a row distribution~$p$ with payoff $\geq v - \varepsilon$
against every column~$c$.\formal{ProbabilityTheory.mwu_approx_minimax} In
particular the minimax value is attained in the limit, recovering the
finite minimax theorem.\formal{ProbabilityTheory.finite_approx_minimax}
\end{thm}

\begin{proof}[Proof idea]
The algorithm keeps a weight per row, multiplied by a factor $e^{\eta}$
whenever that row would have scored against the observed column, then
renormalized to a distribution. A potential argument bounds the gap
between the algorithm's cumulative payoff and the best fixed row by
$O(\sqrt{\log|R| / T})$ after $T$ rounds; taking $T$ large enough drives
the gap below~$\varepsilon$. Those same weights, read as a mixed strategy,
certify the minimax value: the boosting-as-a-game and regret-as-minimax
correspondence, formalized end to end.
\end{proof}

This realized analogy has a transformer face. The multiplicative-weights
distribution is a softmax over accumulated scores, exactly the operation
a softmax-attention head performs over its key scores. Attention is, in
this precise sense, one step of the Hedge update, and the head's induced
concept class is
well-behaved.\formal{TLT.Bridge.softAttention_wellBehaved} The
online$\,\leftrightarrow\,$game analogy thus reaches into the architecture
of \Cref{ch:generalization}: the potential argument that solves a matrix
game is the same normalization that sits inside every attention layer.

\begin{remark}[Why the realized case matters]
\label{rem:realized-matters}
A taxonomy of obstructions would be vacuous if every analogy were blocked.
\Cref{thm:mwu-minimax} shows the frontier of \Cref{sec:missing-equiv} is
worth pushing: an analogy that looked like a metaphor (``learning against an
adversary is like playing a game'') is a theorem, with a softmax at its
center. The obstruction types classify only the analogies not yet discharged.
\end{remark}


% ================================================================
% OPEN PROBLEMS
% ================================================================

\section{Open Problems}
\label{sec:analogies-open}

The open problems here come in two varieties. The first is to discharge a
parallel into a theorem: find the proof that converts a blue dashed edge
into a verified bridge. The second is harder: prove that no such theorem
can exist, that a specific type mismatch is permanently blocked and not
merely waiting for the right idea. Each problem below is tagged a
\emph{known unknown} (KU), sharply posed but unresolved, or an
\emph{unknown known} (UK), where the right bridging object has not yet
been named.

\begin{openprob}[The sample compression conjecture]
\label{op:compression-vc-analogy}
Does every class of VC dimension~$d$ admit a compression scheme of
size~$O(d)$, discharging the compression$\,\leftrightarrow\,$VC analogy
into an equivalence?  The exponential-size direction is machine-checked
(\leanname{fundamental_vc_compression}); the $O(d)$ conjecture is the most
prominent missing witness.  A known unknown~\cite{MoranYehudayoff2016}.
\end{openprob}

\begin{openprob}[MML versus MDL]
\label{op:mml-mdl}
Is there a formal theorem relating Minimum Message Length and Minimum
Description Length: an equivalence on a restricted class of priors, or a
separation?  Both coding objects are formalizable (\leanname{MDLPrinciple},
\leanname{KolmogorovComplexity}); neither direction has a witness.  A known
unknown.
\end{openprob}

\begin{openprob}[Completing the information--PAC-Bayes bridge]
\label{op:itb-pacbayes-bridge}
Hellstr\"om et al.\ partially bridge information-theoretic and PAC-Bayes
bounds via change-of-measure.  Can the partial bridge be completed to a
theorem that recovers each bound from the other with explicit constants,
turning a ``type mismatch'' into a realized analogy?  An unknown
known~\cite{HellstromDurisi2020}.
\end{openprob}

\begin{openprob}[Formalizing the obstruction taxonomy]
\label{op:formalize-taxonomy}
Can ``type mismatch,'' ``one-way theorem,'' and the other obstruction
types be made precise (categorically, say), so that the classification of
the 32 edges becomes a theorem rather than an expert judgment?  An unknown
known: the meta-level vocabulary has no formal definition yet.
\end{openprob}

\begin{openprob}[Attention as Hedge, at scale]
\label{op:attention-hedge}
\Cref{thm:mwu-minimax} identifies a single softmax with one Hedge step.
Does the no-regret guarantee survive composition? Does a stacked
attention model inherit a minimax/regret certificate, or does depth break
the correspondence?  A known unknown on the transformer side
(\leanname{mwu_approx_minimax}, \leanname{softAttention_wellBehaved}).
\end{openprob}

\begin{openprob}[Provable permanence of a type mismatch]
\label{op:type-mismatch-impossibility}
For a specific pair (algorithmic stability versus PAC-Bayes, say), is
there a proof that \emph{no} bridging theorem exists, or is the
``permanence'' of \Cref{sec:type-mismatch} merely the current absence of a
bridge?  A known unknown: the chapter conjectures permanence but proves
none.
\end{openprob}

\begin{openprob}[Where Rademacher beats VC]
\label{op:rademacher-vc-dist}
Characterize the distributions under which the empirical Rademacher
complexity is exponentially smaller than the VC bound.  A sharp
distributional characterization would convert the
Rademacher$\,\leftrightarrow\,$VC ``missing witness'' into a precise
separation theorem.  A known unknown.
\end{openprob}

\begin{openprob}[Ordinal VC versus mind-change ordinals]
\label{op:ordinal-vc-mindchange}
The two transfinite measures are linked here only by a proof-method
mismatch.  Is there a direct theorem relating the ordinal VC dimension
(\leanname{vcdim_to_ordinal_vcdim}) to the mind-change ordinal of
\Cref{ch:ordinals}?  An unknown known~\cite{MartinSharmaStephan2006}.
\end{openprob}

\begin{openprob}[The active--passive gap]
\label{op:active-passive}
Characterize exactly the classes for which label complexity is
exponentially smaller than sample complexity.  A structural
characterization would turn the label$\,\leftrightarrow\,$sample
``missing witness'' into a theorem about which classes benefit from active
querying.  A known unknown.
\end{openprob}

\begin{openprob}[Minimax for continuous attention]
\label{op:minimax-attention}
The verified minimax theorem (\leanname{finite_approx_minimax}) is finite.
Does it lift to the continuous game a softmax-attention head plays over an
unbounded key set, giving a verified value for attention as an adversarial
prediction game?  An unknown known on the transformer side
(\leanname{softAttention_wellBehaved}).
\end{openprob}

\subsection*{Counterfactual exercises}

The figure's central claim is that analogy-failure is taxonomic: thirty-two
independent failures collapse to six types. The exercises that follow probe
whether you can read that claim off the figure itself, using only the two
codes each edge carries. The color and line style name the obstruction type
(why the analogy is not a theorem). The filled dot marks the five edges
where a verified bridge already exists. The exercises work these codes
against each other, especially the cases where both codes apply to the same
edge and appear to conflict.

\begin{enumerate}[leftmargin=*,itemsep=3pt]
\item Count colors and ignore node identities. Thirty-two analogies fail across
only six obstruction types, with type mismatch the plurality (twelve), so ``why
does this analogy fail?'' has a finite-codomain answer rather than a
case-by-case story. Argue that this taxonomic structure is the figure's central
claim.

\item Read only the edges carrying a filled dot. Exactly five of the
thirty-two edges carry a verified bridge, and they cluster on the
complexity-measure side of the graph. Explain why this small dotted subgraph is
the machine-checked fragment while the rest is conceptual analysis or open
conjecture.

\item The pair \nde{mistake\_bound}--\nde{littlestone\_dimension} \emph{is} a
verified equivalence (\leanname{optimal_mistake_bound_eq_ldim}, hence the dot),
yet the edge is blue. Show that the ``missing witness'' names a distinct fact:
as complexity measures in their own right the two diverge when extended to
bandit or partial-monitoring feedback (\Cref{sec:missing-equiv}). Why does the
diagram need both codes here?

\item On \nde{compression\_scheme}--\nde{vc\_dimension}, the dot marks
\leanname{fundamental_vc_compression} (finite VC dimension implies \emph{some}
finite compression scheme, exponential in $d$), but the edge stays blue dashed
because the $O(d)$ sample compression conjecture is open
(\Cref{op:compression-vc-analogy}). Explain how one edge can carry a verified
direction and an open equivalence at once.

\item The edges \nde{rademacher\_complexity}--\nde{vc\_dimension} and
\nde{label\_complexity}--\nde{sample\_complexity} are blue but carry no dot.
These are \emph{established} non-equivalences (Rademacher complexity can fall
far below the VC bound; label complexity can be exponentially smaller than
sample complexity, \Cref{op:active-passive}), not open conjectures. Explain how
the absence of a dot plus the prose distinguishes a known separation from a
live conjecture inside the same color class.

\item Both \nde{covering\_number}--\nde{vc\_dimension} and
\nde{covering\_number}--\nde{rademacher\_complexity} are green ``one-way
theorem'' edges, but only the first carries a dot
(\leanname{sauer_shelah_exp_bound}); the second is Dudley's entropy integral,
stated classically (\Cref{sec:one-way}). Show how the evidence code splits one
obstruction type into a verified bridge and a literature bridge.

\item On \nde{sample\_complexity}--\nde{vc\_dimension} the dot marks
\leanname{fundamental_theorem}, the equivalence between finite VC dimension and
finite sample complexity. The analogy is usually stated as an identity
(``sample complexity just \emph{is} VC dimension''), which erases the two
directions proved by different arguments (\Cref{sec:one-way}). Explain why the
dot guards a two-directional theorem against being read as a definition.

\item Ask whether the six obstruction types are exhaustive: must every
cross-paradigm analogy fail for one of exactly these reasons? A completeness
theorem would require defining the types formally
(\Cref{op:formalize-taxonomy}) and proving the permanence of a single mismatch
(\Cref{op:type-mismatch-impossibility}). Until both are settled, the counts six
and five are a census, not an invariant of the field.
\end{enumerate}
